\documentclass[twocolumn]{aastex631}
\usepackage{amsmath}
\usepackage{booktabs}
\shorttitle{SDSS optical AGN/sSFR matched-control pilot}
\shortauthors{NebulaMind Research Autopilot}
\begin{document}

\title{Optical AGN Hosts and Catalog Specific Star Formation in SDSS DR17: a Selection-Aware Matched-Control Pilot}
\author{NebulaMind Research Autopilot}
\affiliation{NebulaMind Astrophysics Collaboration, San Francisco, CA 94107, USA}
\correspondingauthor{NebulaMind Research Autopilot}
\email{autopilot@nebulamind.ai}

\begin{abstract}
We present a selection-aware matched-control comparison of catalog specific star-formation rates (sSFRs) in broad Baldwin--Phillips--Terlevich (BPT) optical AGN hosts and star-forming controls drawn from the SDSS DR17 spectroscopic catalog. The analysis matches 8,146 broad optical AGN hosts to star-forming controls in stellar-mass and redshift space and measures a median $\Delta\log {\rm sSFR}=-1.309$ dex; at S/N$\geq 10$, the corresponding matched offset is $-0.744$ dex. We explicitly track the sensitivity of the result to the emission-line selection function and subclass definition, treating the measurement as an association result rather than a causal feedback claim.
\end{abstract}

\keywords{galaxies: evolution --- galaxies: active --- galaxies: star formation --- surveys --- methods: data analysis}
\software{Astropy, SciPy, NumPy, Matplotlib, pandas}

\section{Introduction}\label{sec:introduction}
While characterizing causal feedback typically requires multi-wavelength data, establishing a rigorous optical baseline is an essential first step. Here, the optical denominator denotes the SDSS DR17 emission-line parent sample used as the baseline population for downstream comparisons, not a measurement of any physical process by itself. Here we present a selection-aware matched-control pilot analyzing catalog specific star-formation rates (sSFRs) in broad Baldwin--Phillips--Terlevich (BPT) optical AGN hosts and star-forming controls drawn from the SDSS DR17 spectroscopic catalog. Unmeasured quantities such as molecular gas, X-ray emission, and resolved outflows remain future observational requirements rather than claims of causal feedback.


\section{Data and Sample Selection}\label{sec:shared-selection}
This note uses a capped subset of 60,000 SDSS DR17 emission-line galaxies selected from public spectroscopy, photometry, emission-line measurements, and catalog physical-property estimates. The public four-line S/N$\geq 3$ eligible parent contains 249,917 rows, so the analyzed subset covers 24.0\% of that parent. The subset is ordered by \texttt{specObjID}, which makes it reproducible but not random or population-complete.

\begin{deluxetable*}{lrrr}
\tabletypesize{\scriptsize}
\tablecaption{Shared SDSS DR17 selection cascade used before paper-specific quantities.\label{tab:selection-cascade}}
\tablehead{\colhead{Selection stage} & \colhead{Public DR17 rows} & \colhead{Cached rows} & \colhead{Retention vs. spectro-z parent (fraction)}}
\startdata
SpecObj GALAXY, 0.02<z<0.12 & 501,060 & -- & 1.000 \\
plus galSpecInfo/PhotoObj/galSpecExtra and mass/sSFR bounds & 416,554 & -- & 0.831 \\
plus galSpecLine join & 416,554 & -- & 0.831 \\
four BPT lines positive with positive errors & 373,445 & 60,000 & 0.745 \\
four BPT lines S/N$\geq 3$ & 249,917 & 60,000 & 0.499 \\
four BPT lines S/N$\geq 5$ & 176,523 & 42,446 & 0.352 \\
four BPT lines S/N$\geq 10$ & 91,768 & 22,311 & 0.183 \\
\enddata
\tablecomments{Counts are read-only public SDSS DR17 count queries plus the cached local CSV. Cached rows are shown only where the cache applies.}
\end{deluxetable*}

The four-line requirement is strongly selection dependent. In the public counts, S/N$\geq 3$ keeps 33.6\% of the $-12<\log {\rm sSFR}<-11$ parent bin but 94.9\% of the $-10<\log {\rm sSFR}<-9.5$ bin. Therefore every incidence, quenching, density, gas-denominator, or target-vector statement in this analysis is conditional on the four-line emission-line selection.

Local subset versus public catalog marginal checks found no redshift, stellar-mass, or sSFR bin with a subset-minus-public fraction difference above 5 percentage points. The largest absolute differences were 2.03 percentage points in redshift, -1.63 percentage points in stellar mass, and -0.58 percentage points in sSFR. This is a representativeness diagnostic only; it does not make the capped cache random or complete.


\section{Measurements}\label{sec:measurements}
The row-level measurements include redshift, stellar mass, catalog specific star-formation rate, model colors, and four optical line fluxes/errors. BPT labels and all derived denominators are recomputed locally from the cached table. Catalog SFR/sSFR values are treated as low-redshift SDSS physical-property estimates rather than direct resolved gas or feedback measurements \citep{sdssdr17,brinchmann2004,york2000}.
Unless otherwise noted, quoted fraction uncertainties are binomial counting uncertainties from the stated sample sizes, and bracketed intervals are bootstrap confidence intervals.


\section{Flagship integrated result: optical AGN and catalog sSFR}\label{sec:rp1-result}
BPT classes are computed from H$\alpha$, H$\beta$, [O~III]$\lambda5007$, and [N~II]$\lambda6584$ line ratios using the standard Baldwin--Phillips--Terlevich diagram and Kauffmann/Kewley demarcations \citep{baldwin1981,kewley2001,kauffmann2003bpt,kewley2006}. The cached analysis table contains 39,553 star-forming galaxies, 12,234 intermediate/composite objects, 8,146 broad optical AGN, and 67 unclassified objects.

The preferred estimator matches every broad optical AGN host to the nearest star-forming control in standardized $(\log M_\star,z)$ space, with replacement. This is an association design; controls are not matched in morphology, halo mass, gas mass, aperture scale, AGN luminosity, or duty-cycle phase.

Our comparison of broad BPT optical AGN hosts versus star-forming controls at S/N$\geq 3$ matches $N=8,146$ pairs, measuring a median specific star formation rate offset of $\Delta\log {\rm sSFR}=-1.309$ dex with a 95\% bootstrap interval of $[-1.334,-1.282]$ dex. Applying a moderate mass-redshift caliper of $|\Delta\log M_\star|\leq0.05$ and $|\Delta z|\leq0.002$ retains $N=7,867$ pairs (96.6\% target coverage) and yields a median offset of $-1.318$ dex. A deterministic matching run without replacement yields $N=7,419$ matched pairs and a median offset of $-1.446$ dex, though with poorer stellar-mass balance. Raising the line-S/N threshold to 10 leaves $N=1,530$ matched pairs and reduces the median offset to $-0.744$ dex, and a narrower [N II] Seyfert-like proxy yields $N=2,114$ pairs and a median offset of $-0.763$ dex.


\begin{figure*}
\centering
\includegraphics[width=0.73\textwidth]{../figures/fig-bpt.pdf}
\caption{BPT line-ratio diagram for the cached SDSS DR17 optical emission-line subset used in the RP-1 analysis. This figure documents the optical selection and classification boundary; it does not by itself identify causal AGN feedback.}
\label{fig:bpt}
\end{figure*}

\begin{figure*}
\centering
\includegraphics[width=0.86\textwidth]{../figures/fig-matched-offsets.pdf}
\caption{Matched-pair catalog-sSFR offsets for broad BPT optical AGN hosts minus nearest star-forming controls in stellar-mass--redshift space. The large negative offset is robust within the optical emission-line subset but remains selection- and subclass-dependent.}
\label{fig:offsets}
\end{figure*}


\section{Data Availability}\label{sec:data-avail}
The SDSS DR17 spectroscopy, photometry, emission-line measurements, and catalog quantities used here are publicly available through the SDSS data releases and associated catalog papers cited below. A local subset and manifest are retained in the project repository for reproducibility and are available from the corresponding author upon reasonable request.

\section{Conclusion}\label{sec:conclusion}
In the capped SDSS DR17 emission-line subset, broad BPT optical AGN hosts show a median sSFR offset of $-1.309$ dex relative to mass--redshift matched controls, with a 95\% bootstrap interval of $[-1.334,-1.282]$ dex. Although the offset amplitude is highly dependent on the emission-line selection function (decreasing to $-0.744$ dex at S/N$\geq 10$), the interval remains securely negative. These measurements establish a robust optical association baseline, which will require future molecular gas or direct outflow kinematics follow-up to isolate any causal AGN quenching mechanisms.

\acknowledgments
We thank the SDSS collaboration. This manuscript uses public SDSS DR17 data only.


\begin{thebibliography}{}
\bibitem[Abdurro'uf et al.(2022)]{sdssdr17} Abdurro'uf, Accetta, K., Aerts, C., et al. 2022, ApJS, 259, 35
\bibitem[Baldwin et al.(1981)]{baldwin1981} Baldwin, J.~A., Phillips, M.~M., \& Terlevich, R. 1981, PASP, 93, 5
\bibitem[Brinchmann et al.(2004)]{brinchmann2004} Brinchmann, J., Charlot, S., White, S.~D.~M., et al. 2004, MNRAS, 351, 1151
\bibitem[Kauffmann et al.(2003a)]{kauffmann2003bpt} Kauffmann, G., Heckman, T.~M., Tremonti, C., et al. 2003a, MNRAS, 346, 1055
\bibitem[Kewley et al.(2001)]{kewley2001} Kewley, L.~J., Dopita, M.~A., Sutherland, R.~S., Heisler, C.~A., \& Trevena, J. 2001, ApJ, 556, 121
\bibitem[Kewley et al.(2006)]{kewley2006} Kewley, L.~J., Groves, B., Kauffmann, G., \& Heckman, T. 2006, MNRAS, 372, 961
\bibitem[York et al.(2000)]{york2000} York, D.~G., Adelman, J., Anderson, J.~E., Jr., et al. 2000, AJ, 120, 1579

\end{thebibliography}

\end{document}
